\section{Results and the precise logical claim}
\label{sec:results}

There is no single useful meaning of ``simplest'' for these equations.
Number of unknowns, quantified domain, total polynomial degree, arithmetic
circuit length, and coefficient height are different resources.  This report
gives several improvements over the supplied collection, while preserving
its calibrated blank-machine construction.  The most substantial new
integer-domain result is a polynomial in \textbf{24 unrestricted integers},
of exact total degree \textbf{1936}, with a \textbf{167-operation} evaluator.
Its three-equation version has maximum degree 1920 and uses 162 operations.

For strictly positive integer unknowns, the report gives three new
twelve-variable equations with degree/operation pairs
\[
 (1069,123),\qquad (1073,121),\qquad (1081,120).
\]
An exact refactoring of the existing degree-1353 twelve-variable polynomial
also reduces its evaluator from 116 to 113 operations.  A separate
coefficient improvement reduces the radix exponent from 9,765,625 to
6,364,399, and supplies all four resulting fixed integers as actual binary
payloads.  These are complementary improvements, not a claim of a globally
minimal Diophantine representation.

\subsection{Comparison with the supplied snapshot}

All comparisons refer to the repository snapshot
\texttt{ae3007ca4467011e0e051c84acb4a1dd13971b0e}, rather than a moving
\texttt{main}.  The collection's own reconciliation identifies packages
16--23 as the current wave and records the historical alternatives
\cite{collection}.  Table~\ref{tab:comparison} compares like domains and
the particular resource being improved.

\begin{table}[htbp]
\centering\small
\begin{tabular}{p{2.31in}rrr}
\toprule
Single equation & Unknowns & Degree & Operations\\
\midrule
Previous integer, least count and degree (19) & $26\Z$ &2465&192\\
Previous integer, fewest ops at that count (21) &$26\Z$&2657&168\\
\textbf{New signed separating norm} &$\mathbf{24\Z}$&\textbf{1936}&\textbf{167}\\
\midrule
Previous positive, least degree at 12 (20)&12&1145&123\\
New tightest separating weight&12&1069&123\\
New intermediate weight&12&1073&121\\
New simpler weight&12&1081&120\\
\midrule
Previous positive, fewest ops at 12 (21)&12&1353&116\\
Same polynomial, refactored here&12&1353&113\\
\midrule
Previous monic merger (22)&13&319&91\\
Same polynomial, refactored here&13&319&89\\
\bottomrule
\end{tabular}
\caption{Positive integers unless marked $\Z$.  The integer comparison is
with the former minimum-variable family, not with every larger-variable
equation in the collection.  Source sizes: \cite{repo19,repo20,repo21,repo22}.}
\label{tab:comparison}
\end{table}

The new 24-integer equation improves both recorded minima of the former
26-integer family: its degree is below 2465 and its operation count is
below 168.  Larger-variable integer equations still have much lower
degrees; for example package 23 supplies 27 integers at degree 172
\cite{repo23}.  A lower degree is not being claimed across that tradeoff.

Two existing constructions also retain their particular advantages.
Package 18 gives nine natural-number unknowns at degree
$161844\cdot5^{33}+484$, with 20,531 operations; translating every unknown
by one gives nine positive unknowns at the same degree
\cite{repo18}.  The 15-positive-variable polynomial $P_{67}$ remains a
67-operation equation, of degree 155.  The retained-variable system from
package 22 has 16 positive unknowns, seven equations of maximum degree 24,
and 59 operations \cite{baseline67,retained59}.  Both of these latter
circuits are included with the new complete coefficient tuple.  Their
algebra and operation counts are unchanged.  Thus this report does not
claim fewer than nine positive unknowns, fewer than 67 operations for a
single equation, or fewer than 59 operations for a system.

\subsection{The theorem is an equivalence inside ZFC}

Fix the usual arithmetization of finite proofs in the chosen presentation
of ZFC, and let
\[
 \operatorname{Con}(\mathrm{ZFC})\equiv
 \neg\exists p\,\operatorname{Prf}_{\mathrm{ZFC}}
                         (p,\ulcorner\bot\urcorner).
\]
Here $\operatorname{Prf}_{\mathrm{ZFC}}$ is the finite syntactic proof
predicate.  The repository's proof predicate and any ordinary effective
presentation are related by the standard proof translations formalizable
in ZFC; no semantic soundness assumption is used in the equivalence.

Let $\boldsymbol\kappa=(\beta,K_0,M_C,M_F)$ be the four particular
integers defined in Section~\ref{sec:coefficients}.  For the 24-variable
polynomial $P_{24}$ explicitly defined in Section~\ref{sec:integer},
the claim is
\begin{equation}\label{eq:main-logical}
 \boxed{\quad\mathrm{ZFC}\ \vdash\
 \bigl(\neg\exists\boldsymbol x\in\Z^{24}\;
       P_{24}(\boldsymbol x;\boldsymbol\kappa)=0\bigr)
 \ \longleftrightarrow\ \operatorname{Con}(\mathrm{ZFC}).\quad}
\end{equation}
The analogous statements hold over $\Z_{>0}^{12}$ for each displayed
twelve-variable polynomial and over the corresponding domain for every
listed system.  Integer quantifiers here are an arithmetical abbreviation
inside ZFC.  The coefficients are fixed numerals; none is an extra
unknown, an unspecified oracle, or a value defined by waiting for the
machine to halt.

The construction proves the stronger positive equivalence
\begin{equation}\label{eq:existential-chain}
 \exists\boldsymbol x\;P(\boldsymbol x)=0
 \quad\longleftrightarrow\quad
 \operatorname{BlankHalts}(K215)
 \quad\longleftrightarrow\quad
 \exists p\;\operatorname{Prf}_{\mathrm{ZFC}}(p,\ulcorner\bot\urcorner).
\end{equation}
Negating this equivalence gives \eqref{eq:main-logical}.  In particular,
one does not assume $\operatorname{Con}(\mathrm{ZFC})$ while proving the
reduction, and this report does not prove the absence of solutions.
Under the usual hypotheses of the second incompleteness theorem, if ZFC
is consistent it cannot itself prove that absence.  Bare consistency
alone is not used to claim that ZFC cannot prove the opposite existential
sentence; that additional assertion would require a stronger soundness
hypothesis.

\subsection{What is inherited, and what is proved here}

The last equivalence in \eqref{eq:existential-chain} is the calibrated
K215 endpoint.  The supplied Lean source names it
\texttt{blank\_halts\_iff\_zfc\_refutation}, followed by
\texttt{nonhalting\_iff\_zfc\_consistency} \cite{K215}.  A byte-identical
copy and the machine table are included.  This delivery does not rebuild
that Lean development or represent the new number-theoretic reductions
as kernel-checked theorems.

The inherited native-window compiler represents accepting blank-machine
tableaux by Pell equations and fixed masks
\cite{baseline67,normalized,wang}.  Its soundness applies to arbitrary
solutions; its completeness constructs canonical witnesses from every
finite accepting computation.  The signed-domain and odd-scale results
specify which coordinates need positivity and how square and odd
substitutions preserve completeness \cite{normalized,oddscale}.
The source proof files needed for these interfaces accompany the report.

The new proofs supply three separate bridges.  First, rational Pell
saturation recovers an integer auxiliary root after root separation,
allowing a smaller multiplier than the previous separating construction.
Tighter weights then cancel the top homogeneous terms in four of the
sixteen linear conjugate factors.  Second, a signed separation argument,
including the $w=0$ and $T=0$ boundary cases, allows only six coordinates
to retain positivity guards; the resulting lift uses 24 integers.
Third, a concrete Boolean padding controller replaces the old
five-power rounding of the coefficient layout.  All three are ordinary
mathematical arguments using finite arithmetic, elementary algebraic
number theory, Pell identities and finite computation, hence formalizable
in ZFC.

\subsection{Conventions that keep the size claims meaningful}

The degree is the total degree of the actual polynomial after all domain
substitutions and after fixing the four coefficients.  The maximum
residual degree is reported for a system.  A circuit charges every binary
addition, subtraction and multiplication, including multiplication by a
fixed numeral; using a previously computed value and comparing two
values in a system are free.  Squaring is one multiplication.  The
coefficient numerals are input constants, so their bit lengths are
reported separately from circuit length.

Named intermediate expressions are polynomial abbreviations, not
additional existential variables.  Products over formal radicals below
are short descriptions of integer polynomials.  The delivered circuits
evaluate them using only $+,-,\times$, and the report also gives explicit
division-free formulas.  No root extraction or resultant is an uncharged
primitive.  Positive-domain results are always labeled as such; a
twelve-positive-variable equation is not asserted to have the same
soundness over twelve unrestricted integers.

